\section{Introduction}
\label{sec:intro}

Industrial hoisting systems operate under tightly ordered state sequences: normal operation, first-level degradation, second-level degradation, and fault.
Detecting the second-level degradation state early is safety-critical, because it immediately precedes fault occurrence and provides the last actionable window for intervention.
Fixed-horizon future-state prediction---predicting the operating state $\Delta$ steps ahead from a current multivariate sensor window---is a natural formulation for this task.

\textbf{The rare-boundary collapse problem.}
Despite strong aggregate performance, standard time-series classifiers fail on this task in a systematic way.
On the private Hoister dataset used in this work, the second-level degradation state (class~9) appears in only 185 of 37,417 timestamps ($<0.5\%$).
A classifier that never predicts class~9 still achieves high accuracy by modeling the dominant states well.
In our $\Delta{=}1$ experiments, iTransformer~\cite{liu2024itransformer} achieves 81.8\% accuracy but class-9 F1 of 0.000.
DLinear~\cite{zeng2023dlinear}, TimesNet~\cite{wu2023timesnet}, PatchTST~\cite{nie2023patchtst}, and our own SGTONetV4 baseline all share this failure.
We call this the \emph{rare-boundary collapse} failure mode: a classifier collapses on the rare transition state while appearing competent on aggregate metrics.

\textbf{Our approach.}
We propose \textbf{SGTONetV6}, a shift-aware graph and trigger-oriented network that addresses rare-boundary collapse by decoupling two behaviors.
The \emph{base classifier} handles common states conservatively using a patch temporal encoder, graph-aware transition refinement, destination experts, and prototype logits.
The \emph{rare trigger} is a separate branch that computes a scalar rare score from patch-attentive rare context, transition priors, and boundary logits.
At inference, the base prediction is overridden to class~9 only when the rare score exceeds a validation-calibrated threshold \emph{and} the sample satisfies physical boundary and precursor-state constraints:
\begin{equation}
  \yfinal =
  \begin{cases}
    9 & \text{if } \srate \geq \thresh \;\wedge\; \text{boundary\_flag} \;\wedge\; y_t \in \{5, 7\} \\
    \ybase & \text{otherwise.}
  \end{cases}
  \label{eq:override}
\end{equation}
The boundary and precursor constraints encode operational plausibility: class~9 can only follow states~5 or~7, and only near a state-transition boundary.
Both $y_t$ (the predicted current state from the base classifier) and \texttt{boundary\_flag} (derived from the sliding-window label sequence) are observable at inference without oracle access.
This prevents the trigger from firing indiscriminately and producing uncontrolled false alarms.

\textbf{Contributions.}
\begin{enumerate}
  \item We identify and characterize the rare-boundary collapse failure mode in short-horizon hoister future-state prediction, showing that all five tested baselines score class-9 F1 of 0.000 despite high overall accuracy.
  \item We propose SGTONetV6, a dual-mode model with patch-attentive rare context and boundary-constrained inference override.
  \item On the private Hoister case study ($\Delta{=}1$, three file-level splits), SGTONetV6 achieves class-9 F1 of 0.556 and fault macro-F1 of 0.541, while maintaining competitive macro-F1 of 0.623 vs.\ best baseline 0.619.
  \item Ablation over six variants confirms that both the rare override and the boundary constraint are necessary for rare recovery.
\end{enumerate}

\begin{figure}[t]
  \centering
  \includegraphics[width=0.95\linewidth]{figures/fig1_arch.pdf}
  \caption{SGTONetV6 dual-mode architecture. The shared patch temporal encoder feeds two parallel branches. The base classifier (left) handles common states conservatively via graph-aware transition refinement, destination experts, and prototype logits. The rare trigger (right) computes a patch-attentive rare score and overrides the base prediction to class~9 only when the rare score exceeds the calibrated threshold and boundary/precursor constraints are satisfied.}
  \label{fig:arch}
\end{figure}

The rest of this paper is organized as follows.
Section~\ref{sec:related} reviews related work.
Section~\ref{sec:method} describes SGTONetV6.
Section~\ref{sec:experiments} presents experiments and ablations.
Section~\ref{sec:conclusion} concludes with limitations and future directions.
